% GENERATED FILE. Edit canonical lessons, then run npm run generate:books.
% canonical-source-hash: fnv1a64:334945199928324f
% canonical-lessons: KA-C283-nadi, KA-C283-seleta, KA-C283-murita, KA-C283-aparesan, KA-C283-garbhini

\chapter{Pulse, Cramp, Fracture, Operation, Pregnant}
\label{ch:ka-pulse-cramp-fracture-operation-pregnant}
\hlchaptermodality{283}

\begin{chapteropening}
\textbf{By the end of this chapter:} \emph{I can say the pulse, a cramp, a fracture, an operation and pregnant.}

Complete the last lesson of chapter 283: I can say the pulse, a cramp, a fracture, an operation and pregnant.
\end{chapteropening}

\section[\texorpdfstring{\kn{ನಾಡಿ} (nāḍi)}{ನಾಡಿ (nāḍi)}]{\kn{ನಾಡಿ} (nāḍi) — the pulse}

\label{lesson:KA-C283-nadi}



\begin{quote}
Before the new one: say the Kannada for a nurse, then the Kannada for surprise.
\end{quote}

\subsection*{You'll want to know: \kn{ನಾಡಿ}}
\textbf{\kn{ನಾಡಿ}} — \emph{nāḍi} — \textquotedblleft{}the pulse\textquotedblright{}.

\begin{grammarlens}[title={What you've built}]
One more word for this chapter.
\end{grammarlens}

\subsection*{Your turn}
\begin{itemize}
\raggedright
  \item \emph{Say it:} \emph{nāḍi}
  \item \emph{Say it:} \emph{nāḍi}, once more
  \item \emph{Say it:} say \emph{nāḍi}
  \item \emph{Recall:} say the Kannada for a nurse, then the Kannada for surprise, then say \emph{nāḍi} again
\end{itemize}

\begin{tcolorbox}[breakable,colback=teal!4,colframe=teal!35!black,title={Before you move on}]
What does \kn{ನಾಡಿ} mean? (\textquotedblleft{}The pulse\textquotedblright{}.) Say it once more.
\end{tcolorbox}
\section[\texorpdfstring{\kn{ಸೆಳೆತ} (seḷeta)}{ಸೆಳೆತ (seḷeta)}]{\kn{ಸೆಳೆತ} (seḷeta) — a cramp, a spasm}

\label{lesson:KA-C283-seleta}



\begin{quote}
Before the new one: say the Kannada for surprise, then the Kannada for the pulse.
\end{quote}

\subsection*{You'll want to know: \kn{ಸೆಳೆತ}}
\textbf{\kn{ಸೆಳೆತ}} — \emph{seḷeta} — \textquotedblleft{}a cramp, a spasm\textquotedblright{}.

\begin{grammarlens}[title={What you've built}]
One more word for this chapter.
\end{grammarlens}

\subsection*{Your turn}
\begin{itemize}
\raggedright
  \item \emph{Say it:} \emph{seḷeta}
  \item \emph{Say it:} \emph{seḷeta}, once more
  \item \emph{Say it:} say \emph{seḷeta}
  \item \emph{Recall:} say the Kannada for surprise, then the Kannada for the pulse, then say \emph{seḷeta} again
\end{itemize}

\begin{tcolorbox}[breakable,colback=teal!4,colframe=teal!35!black,title={Before you move on}]
What does \kn{ಸೆಳೆತ} mean? (\textquotedblleft{}A cramp, a spasm\textquotedblright{}.) Say it once more.
\end{tcolorbox}
\section[\texorpdfstring{\kn{ಮುರಿತ} (murita)}{ಮುರಿತ (murita)}]{\kn{ಮುರಿತ} (murita) — a fracture}

\label{lesson:KA-C283-murita}



\begin{quote}
Before the new one: say the Kannada for the pulse, then the Kannada for a cramp.
\end{quote}

\subsection*{You'll want to know: \kn{ಮುರಿತ}}
\textbf{\kn{ಮುರಿತ}} — \emph{murita} — \textquotedblleft{}a fracture\textquotedblright{}.

\begin{grammarlens}[title={What you've built}]
One more word for this chapter.
\end{grammarlens}

\subsection*{Your turn}
\begin{itemize}
\raggedright
  \item \emph{Say it:} \emph{murita}
  \item \emph{Say it:} \emph{murita}, once more
  \item \emph{Say it:} say \emph{murita}
  \item \emph{Recall:} say the Kannada for the pulse, then the Kannada for a cramp, then say \emph{murita} again
\end{itemize}

\begin{tcolorbox}[breakable,colback=teal!4,colframe=teal!35!black,title={Before you move on}]
What does \kn{ಮುರಿತ} mean? (\textquotedblleft{}A fracture\textquotedblright{}.) Say it once more.
\end{tcolorbox}
\section[\texorpdfstring{\kn{ಆಪರೇಷನ್} (āpareṣan)}{ಆಪರೇಷನ್ (āpareṣan)}]{\kn{ಆಪರೇಷನ್} (āpareṣan) — an operation, surgery}

\label{lesson:KA-C283-aparesan}



\begin{quote}
Before the new one: say the Kannada for a cramp, then the Kannada for a fracture.
\end{quote}

\subsection*{You'll want to know: \kn{ಆಪರೇಷನ್}}
\textbf{\kn{ಆಪರೇಷನ್}} — \emph{āpareṣan} — \textquotedblleft{}an operation, surgery\textquotedblright{}.

\begin{grammarlens}[title={What you've built}]
One more word for this chapter.
\end{grammarlens}

\subsection*{Your turn}
\begin{itemize}
\raggedright
  \item \emph{Say it:} \emph{āpareṣan}
  \item \emph{Say it:} \emph{āpareṣan}, once more
  \item \emph{Say it:} say \emph{āpareṣan}
  \item \emph{Recall:} say the Kannada for a cramp, then the Kannada for a fracture, then say \emph{āpareṣan} again
\end{itemize}

\begin{tcolorbox}[breakable,colback=teal!4,colframe=teal!35!black,title={Before you move on}]
What does \kn{ಆಪರೇಷನ್} mean? (\textquotedblleft{}An operation, surgery\textquotedblright{}.) Say it once more.
\end{tcolorbox}
\section[\texorpdfstring{\kn{ಗರ್ಭಿಣಿ} (garbhiṇi)}{ಗರ್ಭಿಣಿ (garbhiṇi)}]{\kn{ಗರ್ಭಿಣಿ} (garbhiṇi) — pregnant; a pregnant woman}

\label{lesson:KA-C283-garbhini}



\begin{quote}
Before the new one: say the Kannada for a fracture, then the Kannada for an operation.
\end{quote}

\subsection*{You'll want to know: \kn{ಗರ್ಭಿಣಿ}}
\textbf{\kn{ಗರ್ಭಿಣಿ}} — \emph{garbhiṇi} — \textquotedblleft{}pregnant; a pregnant woman\textquotedblright{}.

\begin{grammarlens}[title={What you've built}]
One more word for this chapter.
\end{grammarlens}

\subsection*{Your turn}
\begin{itemize}
\raggedright
  \item \emph{Say it:} \emph{garbhiṇi}
  \item \emph{Say it:} \emph{garbhiṇi}, once more
  \item \emph{Say it:} say \emph{garbhiṇi}
  \item \emph{Recall:} say the Kannada for a fracture, then the Kannada for an operation, then say \emph{garbhiṇi} again
\end{itemize}

\begin{tcolorbox}[breakable,colback=teal!4,colframe=teal!35!black,title={Before you move on}]
What does \kn{ಗರ್ಭಿಣಿ} mean? (\textquotedblleft{}Pregnant; a pregnant woman\textquotedblright{}.) Say it once more.
\end{tcolorbox}
